\chapter{順動力学・トルク・スケーリング則（Sim04）}\label{ch:dyn}

\section{運動方程式}

導出とスケーリング則の詳細は付録~\ref{app:eom} を参照。


一般化速度 u = (θ̇, q̇) とする。外力なし・全運動量ゼロなら系重心は静止し、 運動エネルギーは T = ½ u\textsuperscript{T}M(q)u（θ に依存しない）。ラグランジュ方程式は次のとおり。

\[
M(q)\dot u + \dot M(q) u - \tfrac12\left[0;\ u^\top \tfrac{\partial M}{\partial q_k} u\right] = [0;\ \tau]
\]

\begin{itemize}
\tightlist
\item
  M の第 1 行は (A, B) そのもので、第 1 式は dL/dt = 0 を表す。
\item
  関節トルク τ は胴体とリンクの間の \textbf{内力} なので、胴体には直接トルクがかからない。
\end{itemize}

\textbf{時間スケーリングの性質（厳密）}: 同じ経路を f 倍速で回すと、τ = M q̈ + (速度の 2 次項) なので \textbf{τ ∝ f²、パワー ∝ f³}。一方、回転角は経路のみで決まり、f に依存しない（第~\ref{ch:theory}章）。

\section{Part A: 検証}

関節トルク τ(t) だけを入力し、開ループで 2 周分（0.5 Hz）を数値積分した。

\begin{longtable}[]{@{}
  >{\raggedright\arraybackslash}p{(\columnwidth - 2\tabcolsep) * \real{0.5000}}
  >{\raggedright\arraybackslash}p{(\columnwidth - 2\tabcolsep) * \real{0.5000}}@{}}
\toprule\noalign{}
\begin{minipage}[b]{\linewidth}\raggedright
検証項目
\end{minipage} & \begin{minipage}[b]{\linewidth}\raggedright
結果
\end{minipage} \\
\midrule\noalign{}
\endhead
\bottomrule\noalign{}
\endlastfoot
関節軌道の追従誤差 & 2.6e-8° \\
2 周後の正味回転 & +13.4086°（運動学モデル Sim03 と一致） \\
角運動量 & max 2e-9 kg·m²/s（保存） \\
エネルギー収支 ΔT − ∫P dt & 9e-4 J（台形積分誤差の範囲） \\
\end{longtable}

→ 運動学モデル（Sim01〜03）と動力学モデルは整合している。

\section{Part B: 人間大（76 kg）}

\FIG{sim04_human_cycle.png}

\begin{longtable}[]{@{}
  >{\raggedright\arraybackslash}p{(\columnwidth - 2\tabcolsep) * \real{0.5000}}
  >{\raggedright\arraybackslash}p{(\columnwidth - 2\tabcolsep) * \real{0.5000}}@{}}
\toprule\noalign{}
\begin{minipage}[b]{\linewidth}\raggedright
項目
\end{minipage} & \begin{minipage}[b]{\linewidth}\raggedright
0.5 Hz で最適ループを回す場合
\end{minipage} \\
\midrule\noalign{}
\endhead
\bottomrule\noalign{}
\endlastfoot
必要トルク peak & 肩 8.2 N·m、股 28.8 N·m \\
パワー peak & 肩 21 W、股 47 W \\
平均正仕事率（回生なし） & 18.5 W \\
持続トルク目安（肩 40、股 120 N·m）での f\_max & \textbf{1.02 Hz} \\
90° 旋回時間（f\_max） & \textbf{13.2 s} \\
\end{longtable}

\begin{itemize}
\tightlist
\item
  人間の筋力で十分実行できる負荷である。90° に十数秒という数値は、宇宙飛行士の 「腕や脚を回して、ゆっくりなら向きを変えられる」という経験と整合する。
\item
  パワーは周期的に正負を繰り返し、1 周あたりの正味仕事はゼロ（内力のみ・摩擦なし）。 \textbf{回生・弾性要素（ばね）を使えば、消費エネルギーは原理的にほぼゼロにできる}。 動物が腱の弾性を使うのと同じ発想で、フィジカル AI の「身体性の活用」の論点になる。
\end{itemize}

\section{Part C: スケーリング則 --- 二乗三乗則が AMBAC を直撃する}

\FIG{sim04_scaling.png}

人間大モデルを長さ s 倍・同密度で相似拡大した（s=10 で約 18 m・76 t）。 アクチュエータは「同じ技術で相似拡大」と仮定し、最大トルク ∝ s³（断面積 s² × レバー s）とした。

\begin{longtable}[]{@{}
  >{\raggedright\arraybackslash}p{(\columnwidth - 4\tabcolsep) * \real{0.3333}}
  >{\raggedright\arraybackslash}p{(\columnwidth - 4\tabcolsep) * \real{0.3333}}
  >{\raggedright\arraybackslash}p{(\columnwidth - 4\tabcolsep) * \real{0.3333}}@{}}
\toprule\noalign{}
\begin{minipage}[b]{\linewidth}\raggedright
量
\end{minipage} & \begin{minipage}[b]{\linewidth}\raggedright
スケーリング（理論）
\end{minipage} & \begin{minipage}[b]{\linewidth}\raggedright
数値で得た指数
\end{minipage} \\
\midrule\noalign{}
\endhead
\bottomrule\noalign{}
\endlastfoot
1 周あたりの回転角 Δθ\_c & s\textsuperscript{0}（不変） & --- \\
同じ周波数で回すのに必要なトルク & s\textsuperscript{5} & 5.00 \\
達成可能な周回周波数 f\_max & s\textsuperscript{-1} & −1.00 \\
90° 旋回時間 & s\textsuperscript{1} & --- \\
AMBAC の角運動量容量 h & s\textsuperscript{4} & 4.00 \\
\end{longtable}

\begin{longtable}[]{@{}
  >{\raggedright\arraybackslash}p{(\columnwidth - 8\tabcolsep) * \real{0.2000}}
  >{\raggedright\arraybackslash}p{(\columnwidth - 8\tabcolsep) * \real{0.2000}}
  >{\raggedright\arraybackslash}p{(\columnwidth - 8\tabcolsep) * \real{0.2000}}
  >{\raggedright\arraybackslash}p{(\columnwidth - 8\tabcolsep) * \real{0.2000}}
  >{\raggedright\arraybackslash}p{(\columnwidth - 8\tabcolsep) * \real{0.2000}}@{}}
\toprule\noalign{}
\begin{minipage}[b]{\linewidth}\raggedright
s
\end{minipage} & \begin{minipage}[b]{\linewidth}\raggedright
全長
\end{minipage} & \begin{minipage}[b]{\linewidth}\raggedright
質量
\end{minipage} & \begin{minipage}[b]{\linewidth}\raggedright
f\_max
\end{minipage} & \begin{minipage}[b]{\linewidth}\raggedright
90° 旋回
\end{minipage} \\
\midrule\noalign{}
\endhead
\bottomrule\noalign{}
\endlastfoot
0.5 & 約 0.9 m & 9.5 kg & 2.0 Hz & 6.6 s \\
1 & 約 1.8 m & 76 kg & 1.0 Hz & 13 s \\
10 & 約 18 m & 76 t & 0.10 Hz & \textbf{132 s} \\
20 & 約 36 m & 608 t & 0.05 Hz & 263 s \\
\end{longtable}

\textbf{18 m 級について} - 人間と同じ機敏さ（0.5 Hz）で回すには、股関節に 2.9 MN·m が必要になる（人間の 10 万倍）。 - 90° を 10 s で旋回するには 20.8 MN·m が必要で、\textbf{相似則で拡大したアクチュエータの 173 倍} にあたる。 - 相似則のままなら、90° の旋回に 2 分以上かかる。

\textbf{考察} - 「回転角はスケール不変」（Sim01）でも、\textbf{速さは 1/s で落ちる}。巨大ロボの AMBAC は遅い。 - 逆に、小さい生物ほど AMBAC が速い。猫が空中で素早く向きを変えられるのは、この則の小さい側にいるため。 - 劇中のように機敏な AMBAC を実現するには、人間の筋肉や既存のモータに比べて \textbf{トルク密度が 2 桁以上高いアクチュエータ} が必要になる。これは物理的な要求として定量化できた。

\section{Part D: スラスタとの比較（18 m 級）}

慣性モーメント約 1.0×10\textsuperscript{6} kg·m²、スラスタのモーメントアーム 5 m、Isp 300 s（化学推進相当）とし、 90° 旋回（静止→静止、加速→減速）に必要な推進剤を求めた。推進剤は 1/T に比例する。

\begin{longtable}[]{@{}
  >{\raggedright\arraybackslash}p{(\columnwidth - 4\tabcolsep) * \real{0.3333}}
  >{\raggedright\arraybackslash}p{(\columnwidth - 4\tabcolsep) * \real{0.3333}}
  >{\raggedright\arraybackslash}p{(\columnwidth - 4\tabcolsep) * \real{0.3333}}@{}}
\toprule\noalign{}
\begin{minipage}[b]{\linewidth}\raggedright
旋回時間 T
\end{minipage} & \begin{minipage}[b]{\linewidth}\raggedright
スラスタ推進剤 / 回
\end{minipage} & \begin{minipage}[b]{\linewidth}\raggedright
AMBAC
\end{minipage} \\
\midrule\noalign{}
\endhead
\bottomrule\noalign{}
\endlastfoot
5 s & 88 kg & 不可能（相似則アクチュエータ） \\
10 s & 44 kg & 不可能（173 倍のトルクが必要） \\
30 s & 15 kg & 不可能 \\
132 s & 3.4 kg & \textbf{推進剤 0 kg}、関節の正仕事 約 2.1 MJ（回生なしの場合） \\
\end{longtable}

\textbf{考察} - ゆっくり回るならスラスタの推進剤も少なくて済む。AMBAC の節約効果が大きくなるのは素早い旋回のときだが、 巨大機では素早い旋回こそ AMBAC が苦手とする領域である。\textbf{トレードオフが逆向きに噛み合っている}。 - ただし、推進剤ゼロという性質は、長期間の運用・推進剤が尽きた状況・排気を出せない状況 （汚染、被探知）では決定的な利点になる。 - 合理的な設計は \textbf{ハイブリッド} になる。素早い旋回はスラスタ（あるいは大トルクの CMG）で行い、 低速の姿勢調整・外乱の吸収・推進剤節約は AMBAC で行う。 「どの時間スケールの運動をどのアクチュエータに配分するか」は、そのまま最適制御・学習の問題になる。

\section{次への論点}

\begin{enumerate}
\def\labelenumi{\arabic{enumi}.}
\tightlist
\item
  \textbf{統合制御（フィジカル AI 編）}: スラスタ・RW/CMG・AMBAC の配分を、推進剤・時間・揺れの多目的で 最適制御と強化学習によって解き、比較する。
\item
  \textbf{弾性・回生}: 直列弾性アクチュエータ（ばね）で周回運動のエネルギーをほぼゼロにできるか。 共振周波数とループ周波数の整合。
\item
  \textbf{3 次元化}: 3 次元ではロール・ピッチ・ヨーの結合（非可換性）により、平面では見えない回転が生じる。
\end{enumerate}
